\section{Limitations}
\label{sec:limitations}
% PLAN (about 0.5 page, one short paragraph per item). Items 1-9 are known now from DEVIATIONS.md, DECISIONS.md and SAP_v2.md; add what the final
% results reveal. State each limitation with its consequence for how the results must be read. Do not soften a limitation that changes a conclusion.
\begin{enumerate}
\item \emph{Synthetic supervision and domain gap.} Training uses rendered or translated projections; the bone image of a real radiograph is never
 observed. Real-image fidelity is measured against dual-energy soft-tissue images, which differ from the radiograph in tone and in high-frequency
 content, so raw and band-limited metrics are both reported and the calibrated metrics are not absolute.
\item \emph{One source of paired real data.} SZCH-X-Rays comes from one partner hospital and one acquisition set-up (dual-exposure DES on a GE Discovery XR656).
\item \emph{The audit is linear and additive.} Lesions are added in image units to a real radiograph; beam hardening, scatter and the tone mapping of
 the real image are not modelled, so $A$ measures the leak of an added signal and not clinical detectability. $A=0$ is reached by doing nothing, and
 $A$ is therefore read next to a suppression gain.
\item \emph{Shared insertion mechanism.} The regulariser and the audit both insert lesions. The confirmatory audit, the real-nodule contrast and the
 downstream subgroup analysis were added to limit this, but they do not remove it.
\item \emph{Development use of SZCH.} All 75 SZCH images were used in development (audit) and 24 for paired decisions; only the 51 held-out pairs are
 clean for fidelity claims, and the SZCH audit is development evidence only. The plan changed once after the SZCH pairs were looked at
 (SAP\_v2.md, DEVIATIONS.md).
\item \emph{Resolution and backbone.} Our models predict at 512 px with tiny encoders (ConvNeXt-V2-T for ours, SwinV2-T for the dose-response);
 the reference model predicts at 1024 px with SwinV2-B. The SSIM gap sits in fine texture: with a band limit of $\sigma=4$ it shrinks from 0.021 to
 0.002, and ours keeps about 9\,\% more high-frequency energy than the dual-energy target in non-bone regions (the reference model 1\,\% less). Whether
 prediction at 1024 px closes the remaining gap was not tested.
\item \emph{Operating point.} The twin weight trades absorption against fidelity and against bone accuracy on exact-truth data (Section~\ref{sec:discussion}).
 Only $w=16$ was evaluated with the full protocol; the dose-response shows that $w=32$ and 64 lower bone-region PSNR by about 0.3\,dB.
\item \emph{Uncertainty.} Conformal calibration treats pixels as exchangeable, so coverage is marginal over pixels, and it is calibrated on synthetic
 data; coverage on real images is reported as measured under shift. The recalibration on real pairs was chosen after the shift was seen and
 used 24 pairs from one hospital. Ensemble members share training data.
\item \emph{Scope.} Training and paired evaluation use adult frontal radiographs; paediatric images enter only the external validation, where no
 paired truth exists; no lateral images and no reader study. Research use only and not a medical device. The weights are non-commercial because
 the training data are.
\item \emph{Overlap maps.} On JSRT the bone overlap of each audit site came from the reference model's bone map, not from the consensus map the
 plan names, so the JSRT interaction is defined by a comparator; on VinDr-RibCXR it comes from radiologist rib masks and shows no interaction.
\item \emph{Downstream and external validation.} Downstream models were trained once per seed with fixed settings, not tuned per input; the
 subgroup analysis by bone overlap is not estimable on TBX11K, where every test box overlaps bone. BDCXR-3257 is a single-centre set with
 image-level labels and no patient identifiers, so repeated images of one child can only be caught by the near-duplicate grouping; its audit
 inserts adult nodules into children's radiographs.
\item \emph{Post-freeze corrections.} Several analysis errors found in an independent check after the first locked run were corrected
 (a resampling error that narrowed interaction intervals, incomplete downstream statistics, and detection metrics that depended on the order of
 tied scores by up to 0.00055); some evaluations ran on cloud GPUs whose floating-point results differ from the local ones by at most
 0.0024\,dB, with the audit sites pinned to the local lung masks. No pre-specified verdict changed; every change is listed in the deviations log.
\item \emph{Registration.} The analysis plans were hash-frozen in the project's version control before testing but not deposited in a public
 registry.
\end{enumerate}
